% ============================================================
%  02_outils.tex — Section 2 : Outils et technologies
% ============================================================
\section{Outils et technologies}

\subsection{Python 3}

Le programme de modélisation et d'orchestration du pipeline est écrit en Python~3. Python a été
choisi pour sa flexibilité, l'écosystème scientifique riche, et la compatibilité directe avec
la bibliothèque PyCSP3. Le solveur de contraintes proprement dit est ACE (Java), appelé en
subprocess depuis Python.

\subsection{PyCSP3 — Modélisation de contraintes}

\begin{definition}[PyCSP3]
  PyCSP3~\cite{pycsp3} est une bibliothèque Python développée par Christophe Lecoutre (CRIL,
  Université d'Artois) qui permet de \textbf{modéliser des problèmes de satisfaction de
  contraintes} au format standard XCSP3~\cite{xcsp3}. Elle fournit une API Python pour
  déclarer des variables, domaines et contraintes, puis génère un fichier XML exploitable
  par n'importe quel solveur compatible XCSP3.
\end{definition}

Principales fonctions utilisées :
\begin{itemize}
  \item \texttt{VarArray(size, dom)} — déclaration d'un tableau de variables avec domaines
  \item \texttt{satisfy(...)} — ajout d'une contrainte au modèle
  \item \texttt{Sum(x) == k} — contrainte de somme
  \item \texttt{scope in table} — contrainte extensionnelle (table de compatibilité)
  \item \texttt{LexIncreasing(x, y)} — $x \leq_{\text{lex}} y$ (brisure de symétrie)
  \item \texttt{compile(filename)} — génération du fichier XCSP3
\end{itemize}

\begin{attention}
  PyCSP3 intercepte \texttt{sys.argv} dès son importation, ce qui entre en conflit avec les
  arguments du script appelant. La solution adoptée consiste à sauvegarder \texttt{sys.argv}
  avant l'import de PyCSP3, puis à le restaurer.
  \begin{lstlisting}
_own_argv = sys.argv[:]
sys.argv = [_own_argv[0]]
from pycsp3 import *
  \end{lstlisting}
\end{attention}

\subsection{ACE — Solveur de contraintes}

\begin{definition}[ACE]
  ACE~\cite{ace} (Adaptive Constraint Engine) est un solveur de contraintes développé par
  Christophe Lecoutre. C'est un programme Java qui lit un fichier XCSP3 et énumère les
  solutions. Il est distribué avec PyCSP3 sous forme de fichier JAR
  (\texttt{ACE-2.5.jar}).
\end{definition}

ACE est appelé en \textbf{subprocess} plutôt que via l'API interne de PyCSP3 (qui retournait
\texttt{UNKNOWN} dans certains cas). La commande utilisée :

\begin{lstlisting}[style=benzai]
java -jar ACE-2.5.jar model.xml -s=all -xe
\end{lstlisting}

Où \texttt{-s=all} demande l'énumération de toutes les solutions et \texttt{-xe} active
l'export XML de chaque solution.

\subsection{XCSP3 — Format standard}

\begin{definition}[XCSP3]
  XCSP3~\cite{xcsp3} est un format XML standardisé pour la représentation de problèmes de
  satisfaction de contraintes. Il est maintenu par le réseau de recherche XCSP
  (\url{https://xcsp.org}) et sert de format d'échange universel entre modélisateurs et
  solveurs.
\end{definition}

\subsection{NetworkX — Graphes en Python}

\begin{definition}[NetworkX]
  NetworkX~\cite{networkx} est une bibliothèque Python pour la création, la manipulation et
  l'étude de graphes. Dans ce projet, elle est utilisée pour :
  \begin{itemize}
    \item Représenter le \textbf{graphe dual} $\GD$ du benzénoïde
    \item Calculer les \textbf{automorphismes} via \texttt{GraphMatcher}
    \item Déterminer les \textbf{degrés} et \textbf{voisinages}
  \end{itemize}
\end{definition}

\subsection{xTB — Validation par chimie quantique semi-empirique}

\begin{definition}[xTB]
  xTB~\cite{xtb} (extended Tight-Binding) est un programme de chimie quantique semi-empirique
  développé par le groupe Grimme (Université de Bonn). La méthode GFN2-xTB~\cite{gfn2xtb}
  traite les électrons de liaison explicitement (contrairement aux champs de force classiques)
  tout en restant rapide. Elle est utilisée dans le pipeline pour optimiser la géométrie des
  solutions reconstruites et vérifier leur planarité.
\end{definition}

\subsection{Benzai — Générateur de benzénoïdes}

\begin{definition}[Benzai]
  Benzai~\cite{benzai} est un logiciel de génération de structures benzénoïdes développé au
  LS2N (Université de Nantes). Il produit des fichiers \texttt{.graph} au format DIMACS étendu,
  décrivant la topologie d'un benzénoïde (sommets = carbones, arêtes = liaisons, faces =
  hexagones). Ces fichiers constituent l'entrée du programme CSP.
\end{definition}
